\section{Implementación serial}

\subsection{Estructura del código}

Siguiendo la organización modular del proyecto, el código se separa en un módulo por algoritmo y un módulo común para el generador de números aleatorios:

\begin{center}
\begin{tabular}{|l|p{8.5cm}|}
\hline
\textbf{Archivo} & \textbf{Contenido} \\
\hline
\texttt{modules/rng.c} & Generador \texttt{xoshiro256**} con mezcla de semillas \texttt{SplitMix64} \\
\texttt{modules/pi\_dartboard.c} & Kernel del dartboard: versión serial y versiones OpenMP \\
\texttt{modules/pi\_buffon.c} & Kernel de Buffon: versión serial y versiones OpenMP \\
\texttt{modules/pi\_common.c} & Despacho entre variantes, estimador y desviación teórica \\
\texttt{src/montecarlo.c} & Programa principal y modo de validación \\
\hline
\end{tabular}
\end{center}

El programa principal recibe el algoritmo, el número de muestras, la cantidad de hilos, la variante y la semilla, y devuelve los resultados en formato clave=valor, legible tanto por una persona como por el script de medición.

\subsection{El generador de números aleatorios}

Es la pieza crítica de todo el trabajo, y merece un apartado propio porque condiciona tanto la corrección como el paralelismo.

Se usa \textbf{xoshiro256**}, un generador de 256 bits con periodo $2^{256}$ y excelente calidad estadística, con dos detalles importantes:

\begin{itemize}
    \item \textbf{Se toman los 53 bits más significativos} para construir el \texttt{double} en $[0,1)$. Los bits bajos de un generador congruencial lineal tienen periodos muy cortos y no sirven.
    \item \textbf{Las semillas de cada hilo se derivan con \texttt{SplitMix64}} a partir de la semilla base y el índice del hilo. \texttt{SplitMix64} es un mezclador cuyo único propósito es dispersar semillas: garantiza que dos hilos con índices consecutivos (0 y 1, 2 y 3, \ldots) partan de estados muy distintos.
\end{itemize}

El segundo punto no es un detalle menor. Con un LCG clásico, dos generadores con semillas que difieren en 1 producen flujos \emph{fuertemente correlacionados}, y la estimación de $\pi$ queda sesgada. Es un error fácil de cometer y que invalida silenciosamente el resultado.

\subsection{Determinismo}

Con una misma semilla, la versión serial produce exactamente el mismo resultado en cada ejecución. Esto es lo que permite usar la corrida serial como referencia de las versiones paralelas.

\subsection{Las versiones paralelas no son idénticas, y está bien}

Las versiones con OpenMP reparten las $N$ muestras entre los hilos y cada hilo usa su propio flujo. El conjunto de muestras que se genera depende por lo tanto del número de hilos: la versión serial y la de 8 hilos no estiman $\pi$ con los mismos datos. \emph{Elegimos} que sea así, por dos razones:

\begin{itemize}
    \item \textbf{Es lo correcto estadísticamente.} Ambas estimaciones son no sesgadas y convergen al mismo valor con la misma tasa de error. Forzar que la versión de 8 hilos genere exactamente los mismos números que la serial obligaría a pre-generar las $N$ muestras en un arreglo, es decir, a pasar de un espacio de trabajo de unos pocos bytes a uno de $O(N)$, que es justo lo que este reto pide optimizar en memoria.
    \item \textbf{Es lo que hace el paralelismo real.} En una aplicación real, dos hilos, dos procesos o dos máquinas no pueden garantizar el mismo flujo de números aleatorios sin coordinarse, y el costo de esa coordinación superaría el beneficio.
\end{itemize}

Por eso la comparación entre variantes se hace por equivalencia \emph{estadística} (sección 6), no por igualdad exacta.

\subsection{Modo de validación}

El programa incluye un modo \texttt{validar} que repite la corrida varias veces con semillas distintas y reporta la media, la desviación, el mínimo y el máximo de $\hat\pi$, la desviación teórica, el \emph{z-score} (cuántas $\sigma$ teóricas se aparta la media de $\pi$) y un veredicto. Ese veredicto es el que se usa como control antes de medir tiempos: si un algoritmo no converge, no tiene sentido cronometrarlo.
